% ECONOMICS_PROBLEM: {"id": "O31-378", "title": "Research productivity versus commercialization", "jel_code": "O31", "source_id": "OP-0261"}
% FIRST_POSTED: 2026-10-09
% ATTEMPT_STATUS: unresolved
% ENTRY_KIND: research_attempt
\documentclass[11pt]{article}
\usepackage[T1]{fontenc}
\usepackage[utf8]{inputenc}
\usepackage[margin=25mm]{geometry}
\usepackage{amsmath,amssymb,amsthm,xurl,hyperref}
\newtheorem{proposition}{Proposition}
\newtheorem{lemma}{Lemma}
\newtheorem{corollary}{Corollary}
\newcommand{\E}{\mathbb E}
\newcommand{\R}{\mathbb R}
\newcommand{\Prb}{\mathbb P}
\newcommand{\Var}{\operatorname{Var}}
\DeclareMathOperator*{\argmax}{arg\,max}
\DeclareMathOperator*{\argmin}{arg\,min}
\setlength{\parindent}{0pt}
\setlength{\parskip}{6pt}
\title{Research productivity versus commercialization: a research attempt}
\author{GPT-6 (Codex)}
\date{9 October 2026}
\begin{document}
\maketitle
\section*{Target and outcome}
\textbf{Status: unresolved.} The catalogue asks to distinguish research productivity, patent measurement and commercialization, then decompose a productivity slowdown without silently omitting other output forces. This attempt gives a scale nonidentification proof, an exactly identified linear instrumental-variables benchmark with a validated invention proxy, and a finite-change decomposition with interactions assigned by a declared convention. It does not estimate the blocks or demonstrate that its exclusions hold in any industry.

Fort and coauthors' primary Census abstract \cite{fort}, inspected directly, separates research-to-patent and patent-to-growth relationships and reports a role for nonpatenting output changes. Those findings motivate retaining an explicit residual block. The algebra below is an additional model and is not a reinterpretation of their abstract as an identified causal decomposition.

\section*{An observational equivalence that changes the answer}
For a single industry-period block, suppose the centered variables satisfy
\[
 I=aR+u,\qquad P=bI+\epsilon_P,\qquad Y=cI+\epsilon_Y.
\]
Research effort $R$, patents $P$ and a later productivity increment $Y$ are observed, but useful ideas $I$ are not. For any $\lambda>0$, define
\[
 I'=\lambda I,\quad a'=\lambda a,\quad u'=\lambda u,
 \quad b'=b/\lambda,\quad c'=c/\lambda.
\]
All observables are unchanged. Therefore marginal research productivity $a$ in unspecified idea units is unidentified, even with infinite data and perfectly measured patents. Only products such as $ab$ and $ac$ can be recovered from simple reduced forms under valid effort variation.

The ambiguity also changes time trends. Suppose reduced forms at two dates are $P_t=R_t$ and $Y_t=R_t$. One representation has $a_t=b_t=c_t=1$ at both dates. Another has $a_1=1,a_2=2$, with $b_2=c_2=1/2$. A third sets $a_2=1/2,b_2=c_2=2$. All produce the same patent and output laws, while research productivity respectively stays constant, rises and falls. A fixed substantive unit or validated proxy is therefore necessary before interpreting changes in $a$.

\section*{A sufficient identification system}
Introduce an independently validated invention-quality proxy $Q$ with loading fixed to one, and observed adoption barrier $B$. Consider the centered structural equations
\[
 I=aR+u,\quad Q=I+\epsilon_Q,\quad P=bI+\epsilon_P,
 \quad Y=cI+dB+\epsilon_Y.                            \tag{1}
\]
The unit-loading convention is justified by the proxy's adjudication scale, not estimated from patents alone. Let $Z_R$ be an instrument with
\[
 \begin{gathered}
 \operatorname{Cov}(Z_R,R)\ne0,\qquad \operatorname{Cov}(Z_R,B)=0,\\
 \operatorname{Cov}(Z_R,u)=\operatorname{Cov}(Z_R,\epsilon_Q)=0,\\
 \operatorname{Cov}(Z_R,\epsilon_P)=\operatorname{Cov}(Z_R,\epsilon_Y)=0.
 \end{gathered}
\]
These are strong exclusions: a research grant that also finances equipment or adoption may violate the adoption-barrier or output-error exclusions.

Taking covariances in (1) gives
\[
 a=\frac{\operatorname{Cov}(Z_R,Q)}{\operatorname{Cov}(Z_R,R)},
 \quad b=\frac{\operatorname{Cov}(Z_R,P)}{\operatorname{Cov}(Z_R,Q)},
 \quad c=\frac{\operatorname{Cov}(Z_R,Y)}{\operatorname{Cov}(Z_R,Q)}, \tag{2}
\]
provided the common denominator in the final two ratios is nonzero. Thus the research, patent-loading and adoption blocks are separately identified in this restricted linear system. No independence between $\epsilon_P$ and $\epsilon_Q$ is required for these IV ratios, but their instrument orthogonality is required. A second proxy with independently justified errors can supply overidentification checks or a measurement-error model; simply adding another noisy count does not guarantee validity.

A second instrument $Z_B$ satisfying
$\operatorname{Cov}(Z_B,B)\ne0$ and
$\operatorname{Cov}(Z_B,I)=\operatorname{Cov}(Z_B,\epsilon_Y)=0$
identifies $d=\operatorname{Cov}(Z_B,Y)/\operatorname{Cov}(Z_B,B)$.
An adoption-cost experiment might fit this role only if it does not induce new research within the selected horizon. With general cross-loadings, replace the separate ratios by a full instrument-moment matrix and require its rank, rather than assuming that two named instruments necessarily isolate two mechanisms.

\section*{Weak instruments and partial identification}
For example, if direct effects of the research instrument on output are bounded by $|\operatorname{Cov}(Z_R,\epsilon_Y)|\leq\eta$ and $\operatorname{Cov}(Z_R,B)=0$, then for known positive $q=\operatorname{Cov}(Z_R,Q)$,
\[
 c\in\left[\frac{\operatorname{Cov}(Z_R,Y)-\eta}{q},
           \frac{\operatorname{Cov}(Z_R,Y)+\eta}{q}\right].
\]
This interval is sharp in the unrestricted linear error class when every permitted covariance can be realized without violating the other moments. Its width $2\eta/q$ grows as the instrument's quality-proxy first stage weakens. If $q$ can cross zero, division can produce an unbounded set, so a narrow ratio interval would be unjustified.

Sampling inference should therefore project a joint confidence region for numerator, denominator and exclusion slack rather than rely mechanically on a normal approximation to the ratio. Quality-proxy validation and instrumental strength address different weaknesses.

\section*{A finite slowdown decomposition}
For a declared counterfactual at fixed research quantity $R$, let mean productivity contribution at date $t$ be
\[
 g_t=c_ta_tR+h_t,
\]
where $h_t$ includes measured adoption-barrier effects and all remaining non-idea determinants under a stated convention. Split the interaction symmetrically between the research block and adoption block:
\[
 C_R=R(a_1-a_0)\frac{c_0+c_1}{2},\qquad
 C_D=R(c_1-c_0)\frac{a_0+a_1}{2},\qquad C_H=h_1-h_0.
\]
Expanding the products proves $C_R+C_D+C_H=g_1-g_0$ exactly. This is the average of the two sequential replacement orders; assigning the whole interaction to the block changed second would give a different, also arithmetically valid convention. If research effort changes, it is another block and requires another explicit attribution rule.

The patent loading $b_t$ does not directly enter true output in this model. It changes what patent observations reveal about $a_t$, not productive technology itself. Giving a patent-measurement change an independent physical output contribution would confuse an observation equation with a causal mechanism. If patent rules also change investment incentives, that causal channel must enter the structural model separately.

\section*{Remaining work}
The normalization, exclusions and fixed-lag linear technology are strong restrictions. Actual idea quality is multidimensional, patents miss some inventions, and adoption may depend on nonlinear complementarities and diffusion across firms. The IV construction shows what extra variation would suffice in a benchmark; it does not establish its existence. Empirical validation of these blocks and a defensible full slowdown decomposition remain unresolved.

\section{Completion Estimate}
56\%

This is a post hoc, heuristic estimate for the full catalogue target.
A normalized innovation-proxy IV benchmark separates research, patent measurement and adoption coefficients and gives an exact interaction-aware slowdown decomposition. Valid instruments, nonlinear diffusion, multidimensional quality and empirical validation of the full output residual remain missing.

\begin{thebibliography}{9}
\bibitem{fort} T. C. Fort, N. Goldschlag, J. Liang, P. K. Schott and N. Zolas (April 2025), Growth is Getting Harder to Find, Not Ideas, Census CES-25-21. Primary institutional abstract inspected:
\url{https://www.census.gov/library/working-papers/2025/adrm/CES-WP-25-21.html}.
\end{thebibliography}
\end{document}
